\chapter{Abdominal Pain During Pregnancy}

\section*{Definition}
Abdominal pain during pregnancy is discomfort or pain between the lower ribs and pelvis in a pregnant person. Mild, brief pain can result from normal stretching, gas, constipation, reflux, musculoskeletal strain, or irregular uterine tightenings, but severe, persistent, localized, or bleeding-associated pain may indicate a pregnancy complication or a non-obstetric emergency.

\section*{When to See a Doctor}
Seek emergency care for severe or worsening pain, fainting, dizziness, shoulder-tip pain, heavy vaginal bleeding, fluid leakage, fever or feeling very unwell, repeated vomiting or inability to keep fluids down, painful urination with fever or flank pain, abdominal trauma, regular contractions or pelvic pressure before 37 weeks, reduced fetal movement later in pregnancy, chest pain, shortness of breath, severe headache, visual changes, or persistent upper-right abdominal pain with swelling or high blood pressure. Contact the maternity or obstetric service promptly if pain is severe or does not improve after 30--60 minutes of rest, even if no other warning sign is present.

\section*{How It Develops}
Pregnancy enlarges and changes the uterus, abdominal organs, blood vessels, and pelvic ligaments. Hormonal effects slow bowel movement and relax smooth muscle. As pregnancy progresses, stretching of the round ligaments and pressure from the uterus can cause pain. Dangerous causes can involve ectopic pregnancy or miscarriage in early pregnancy; placental abruption, preterm labor, or uterine rupture later in pregnancy; and infection, appendicitis, gallbladder disease, kidney disease, or intestinal obstruction at any stage. Preeclampsia and HELLP syndrome can cause upper-abdominal pain, usually after 20 weeks or shortly after delivery.

\section*{Causes}
\begin{itemize}
  \item Round-ligament pain, irregular Braxton Hicks tightenings, constipation, gas, reflux, or musculoskeletal strain
  \item Ectopic pregnancy or miscarriage, particularly in early pregnancy
  \item Placental abruption, preterm labor, or, rarely, uterine rupture; risk is higher with a previous uterine scar or during labor
  \item Urinary tract infection, kidney stone, pyelonephritis, or ovarian torsion
  \item Appendicitis, gallbladder disease, pancreatitis, gastroenteritis, or bowel obstruction
  \item Preeclampsia or HELLP syndrome when upper-right abdominal pain occurs with headache, visual symptoms, swelling, nausea or vomiting, or high blood pressure
\end{itemize}

\section*{Related Symptoms}
\begin{itemize}
  \item Vaginal bleeding, pelvic pressure, regular contractions, unusual discharge, fluid leakage, or reduced fetal movement later in pregnancy
  \item Fever, feeling very unwell, vomiting, diarrhea, urinary burning, frequent urination, flank pain, or jaundice
  \item Headache, visual disturbance, sudden swelling of the face or hands, high blood pressure, or shortness of breath
\end{itemize}

\section*{Medical Evaluation}
\begin{itemize}
  \item History includes gestational age, pain location and onset, severity and pattern, bleeding, contractions, fluid leakage, fetal movement, urinary and gastrointestinal symptoms, fever, trauma, previous abdominal or uterine surgery, and prior pregnancy complications.
  \item Examination assesses vital signs, blood pressure, abdominal tenderness or rigidity, uterine activity, fetal heart rate, hydration, and signs of shock, infection, preeclampsia, or placental complications. Pelvic or speculum examination and cervical assessment may be indicated.
  \item Testing may include urinalysis and urine culture, complete blood count, kidney and liver tests, platelets, lipase, blood group and Rh status, quantitative hCG when early pregnancy location or viability is uncertain, and other pregnancy-related tests as indicated.
  \item Ultrasound is usually the first imaging test for early-pregnancy pain or bleeding and for many obstetric or pelvic causes. Fetal monitoring and obstetric ultrasound may be needed later in pregnancy. MRI can be used when it will answer the clinical question; CT should not be withheld when necessary for a serious maternal condition, although MRI or ultrasound is preferred when equally effective. Imaging and testing are selected by gestational age, symptoms, examination, and suspected cause.
\end{itemize}

\section*{Treatment Options}
Treatment is cause-specific and may include supervised hydration, observation by an obstetric or medical team, pregnancy-appropriate analgesia, antibiotics for confirmed or strongly suspected infection, treatment of urinary disease, management of preeclampsia, or urgent obstetric, surgical, or neonatal care. Suspected ectopic pregnancy, placental abruption, uterine rupture, appendicitis, or obstruction requires urgent specialist assessment. Imaging and treatment decisions balance maternal and fetal safety without withholding necessary emergency care.

\section*{Self-Care and Home Management}
Mild pain that clearly improves with rest or position change and has no warning signs may be monitored briefly. Rest, change position, and maintain ordinary hydration as tolerated; gentle activity is acceptable unless a clinician has advised otherwise. Use only medicines recommended as safe in pregnancy by a clinician or pharmacist. Do not use laxatives or pain medicines to mask severe or persistent pain, and do not delay obstetric assessment. Contact the maternity or obstetric service if pain does not settle promptly or if any warning sign develops.

\section*{References}
\begin{thebibliography}{99}
\bibitem{abdominal_pain_pregnancy:ref1} American College of Obstetricians and Gynecologists. Diagnostic Imaging During Pregnancy and Lactation. ACOG Committee Opinion No. 723, reaffirmed 2021.
\bibitem{abdominal_pain_pregnancy:ref2} American College of Obstetricians and Gynecologists. Gestational Hypertension and Preeclampsia. Practice Bulletin No. 222. \textit{Obstet Gynecol}. 2020;135:e237--e260.
\bibitem{abdominal_pain_pregnancy:ref3} Cunningham FG, et al. \textit{Williams Obstetrics}. 26th ed. McGraw-Hill; 2022.
\bibitem{abdominal_pain_pregnancy:ref4} American College of Obstetricians and Gynecologists. Preterm Labor and Birth. Patient guidance, current online resource.
\bibitem{abdominal_pain_pregnancy:ref5} American College of Obstetricians and Gynecologists. Urgent Maternal Warning Signs and Symptoms. Current online resource.
\bibitem{abdominal_pain_pregnancy:ref6} National Health Service. Stomach Pain in Pregnancy. Reviewed 2024.
\end{thebibliography}
